\chapter{Conclusion and Recommendations}

\section{Conclusion}
Most female university students in this study (80.3\%) would accept HPV vaccination, but only a small proportion had been vaccinated. Knowledge was poor in two thirds of students, including most healthcare students, while attitudes were strongly favourable. The main reason for not being vaccinated was not knowing about the vaccine, and the main barriers were lack of awareness, cost, absence of a provider's recommendation and limited availability. Having heard of HPV and a more positive attitude were associated with acceptance, whereas academic stream was not. The results suggest that the main task is to make the vaccine known and reachable to young women, rather than to overcome opposition to it.

\section{Recommendations}
\begin{enumerate}[leftmargin=*,itemsep=3pt]
\item Run awareness campaigns on campuses that explain what HPV is, who can be vaccinated, the protection offered and its limits, and where the vaccine can be obtained.
\item Include HPV, cervical cancer prevention and screening in the curriculum for students of all faculties, and strengthen vaccine counselling content for healthcare students.
\item Encourage doctors and campus health services to recommend HPV vaccination routinely to young women.
\item Consider catch-up vaccination for young women at free or reduced cost, available at campus health centres.
\item Provide accurate, easily shared material on social media and counter misinformation.
\item Carry out larger, multi-institution studies with random sampling, including rural and private institutions, and follow-up studies to see whether intention turns into vaccination.
\end{enumerate}
